% Exact source fragments, not a standalone TeX document.
% B40.add-vector statement; src/vs/vs2.tex; source lines 492-499
\begin{lemma}  \label{lm:AddVecLiSetIsLiIffVecNotInSpan}
%<*lm:AddVecLiSetIsLiIffVecNotInSpan>
Suppose that $S$ is linearly independent and that $\vec{v}\notin S$.
Then
the set $S\union\set{\vec{v}}$ is linearly independent if and only if 
$\vec{v}\notin\spanof{S}$.   
%</lm:AddVecLiSetIsLiIffVecNotInSpan>
\end{lemma}

% B40.add-vector proof; src/vs/vs2.tex; source lines 501-525
\begin{proof}
%<*pf:AddVecLiSetIsLiIffVecNotInSpan0>
% Assume that $S$ is linearly independent and that $\vec{v}\notin S$.
We will show that $S\union\set{\vec{v}}$ is not linearly independent 
if and only if $\vec{v}\in\spanof{S}$.

Suppose first that $\vec{v}\in\spanof{S}$.
Express $\vec{v}$ as a combination $\vec{v}=c_1\vec{s}_1+\cdots+c_n\vec{s}_n$.
Rewrite that $\zero=c_1\vec{s}_1+\cdots+c_n\vec{s}_n-1\cdot\vec{v}$.
Since $\vec{v}\notin S$, it does not equal any of the $\vec{s}_i$ so this is 
a nontrivial linear dependence among the elements of $S\union\set{\vec{v}}$.
Thus that set is not linearly independent.
%</pf:AddVecLiSetIsLiIffVecNotInSpan0>

%<*pf:AddVecLiSetIsLiIffVecNotInSpan1>
Now suppose that $S\union\set{\vec{v}}$ is not linearly independent and
consider a nontrivial dependence among its members
$\zero=c_1\vec{s}_1+\cdots+c_n\vec{s}_n+c_{n+1}\cdot\vec{v}$.
If $c_{n+1}=0$ then that is a dependence among the elements of~$S$, but 
we are assuming that~$S$ is independent, so $c_{n+1}\neq 0$.
Rewrite the equation  
as $\vec{v}=(c_1/c_{n+1})\vec{s}_1+\cdots+(c_n/c_{n+1})\vec{s}_n$
to get $\vec{v}\in\spanof{S}$
%</pf:AddVecLiSetIsLiIffVecNotInSpan1>
\end{proof}

% B40.exchange statement; src/vs/vs3.tex; source lines 1657-1667
\begin{lemma}[Exchange Lemma] \label{lm:ExchangeLemma}
%<*lm:ExchangeLemma>
Assume that 
\( B=\sequence{\vec{\beta}_1,\dots,\vec{\beta}_n} \) is a basis for a
vector space, and that for the vector \( \vec{v} \)
the relationship \( \vec{v}=\lincombo{c}{\vec{\beta}} \)
has \( c_i\neq 0 \).
Then exchanging \( \vec{\beta}_i \) for \( \vec{v} \) yields another
basis for the space.
%</lm:ExchangeLemma>
\end{lemma}

% B40.exchange proof; src/vs/vs3.tex; source lines 1669-1721
\begin{proof}
%<*pf:ExchangeLemma0>
Call the outcome of the exchange 
\( \hat{B}=\sequence{\vec{\beta}_1,\dots,\vec{\beta}_{i-1},\vec{v},
                       \vec{\beta}_{i+1},\dots,\vec{\beta}_n}  \).

We first show that $\hat{B}$ is linearly independent.
Any relationship 
\( d_1\vec{\beta}_1+\dots+d_i\vec{v}+\dots+d_n\vec{\beta}_n=\zero \)
among the members of $\hat{B}$, after substitution for $\vec{v}$,
\begin{equation*}
  d_1\vec{\beta}_1+\dots
    +d_i\cdot(c_1\vec{\beta}_1+\dots+c_i\vec{\beta}_i+\dots+c_n\vec{\beta}_n)
    +\dots+d_n\vec{\beta}_n
  =\zero
\tag*{($*$)}\end{equation*}
gives a linear relationship among the members of $B$.
The basis $B$ is linearly independent so the coefficient $d_ic_i$ of 
$\vec{\beta}_i$ is zero.
Because we assumed that $c_i$ is nonzero, $d_i=0$.
Using this in equation~$(*)$ gives that all of the other $d$'s are also
zero.
Therefore $\hat{B}$ is linearly independent.
%</pf:ExchangeLemma0>

%<*pf:ExchangeLemma1>
We finish by showing that $\hat{B}$ has the same span as $B$.
Half of this argument, that $\spanof{\smash{\hat{B}}}\subseteq\spanof{B}$, 
is easy; we can write any member 
$d_1\vec{\beta}_1+\dots+d_i\vec{v}+\dots+d_n\vec{\beta}_n$
of $\spanof{\smash{\hat{B}}}$ as
$
  d_1\vec{\beta}_1+\dots
    +d_i\cdot(c_1\vec{\beta}_1+\dots+c_n\vec{\beta}_n)
    +\dots+d_n\vec{\beta}_n
$,
which is a linear combination of linear combinations of members of $B$, and
hence is in $\spanof{B}$.
For the $\spanof{B}\subseteq\spanof{\smash{\hat{B}}}$ half of the argument,
recall that if
$\vec{v}=c_1\vec{\beta}_1+\dots+c_n\vec{\beta}_n$ with $c_i\neq 0$
then we can rearrange the equation to
$\vec{\beta}_i=(-c_1/c_i)\vec{\beta}_1+\dots+(1/c_i)\vec{v}+\dots
   +(-c_n/c_i)\vec{\beta}_n$.
Now, consider any member
$d_1\vec{\beta}_1+\dots+d_i\vec{\beta}_i+\dots+d_n\vec{\beta}_n$
of $\spanof{B}$, substitute for $\vec{\beta}_i$ its expression as a linear
combination of the members of $\hat{B}$, and recognize, 
as in the first half of this argument, that the result is a linear
combination of linear combinations of members of $\hat{B}$, and hence is in 
$\spanof{\smash{\hat{B}}}$.
%</pf:ExchangeLemma1>
\end{proof}

% B40.equal-basis-size statement; src/vs/vs3.tex; source lines 1723-1728
\begin{theorem} \label{th:AllBasesSameSize}
%<*th:AllBasesSameSize>
In any finite-dimensional vector space, all 
bases have the same number of elements.
%</th:AllBasesSameSize>
\end{theorem}

% B40.equal-basis-size proof; src/vs/vs3.tex; source lines 1730-1779
\begin{proof}
%<*pf:AllBasesSameSize0>
Fix a vector space with at least one finite basis.
Choose, from among all of this space's bases,
one \( B=\sequence{\vec{\beta}_1,\dots,\vec{\beta}_n} \) of minimal size.
We will show that any other basis
\( D=\smash{\sequence{\vec{\delta}_1,\vec{\delta}_2,\ldots}} \)
also has the same number of members, $n$.
Because \( B \) has minimal size, \( D \) has no fewer than \( n \) vectors.
We will argue that it cannot have more than \( n \) vectors.
%</pf:AllBasesSameSize0>

%<*pf:AllBasesSameSize1>
The basis \( B \) spans the space and \( \vec{\delta}_1 \) is in the space,
so \( \vec{\delta}_1 \) is a nontrivial linear combination of elements of
\( B \).
By the Exchange Lemma, we can swap \( \vec{\delta}_1 \) for a
vector from \( B \), resulting in a basis \( B_1 \), where one element is
\( \vec{\delta}_1 \) and all of the \( n-1 \) other elements 
are \( \vec{\beta} \)'s.
%</pf:AllBasesSameSize1>

%<*pf:AllBasesSameSize2>
The prior paragraph forms the basis step for an induction argument.
The inductive step starts with a basis \( B_k \) (for \( 1\leq k<n \))
containing \( k \) members of \( D \) and \( n-k \) members of \( B \).
We know that \( D \) has at least \( n \) members so there is a
\( \vec{\delta}_{k+1} \).
Represent it as a linear combination of elements of \( B_k \).
The key point:~in that representation, at least one of the nonzero scalars
must be associated with a \( \vec{\beta}_i \) or else that
representation would be a
nontrivial linear relationship among elements of the linearly independent
set \( D \).
Exchange \( \vec{\delta}_{k+1} \) for \( \vec{\beta}_i \) to get a new basis
\( B_{k+1} \) with one \( \vec{\delta} \) more and one \( \vec{\beta} \)
fewer than the previous basis \( B_k \).
%</pf:AllBasesSameSize2>

%<*pf:AllBasesSameSize3>
Repeat that until no \( \vec{\beta} \)'s remain, so
that \( B_n \) contains  
$\vec{\delta}_1,\dots,\vec{\delta}_n$.
Now, \( D \) cannot have more than these \( n \) vectors because
any \( \vec{\delta}_{n+1} \) that remains would be in the span of
\( B_n \) (since it is a basis) 
and hence would be a linear combination of the other $\vec{\delta}$'s, 
contradicting that $D$ is linearly independent.
%</pf:AllBasesSameSize3>
\end{proof}

% B40.dimension-bound statement; src/vs/vs3.tex; source lines 1818-1823
\begin{corollary} \label{cor:NoLiSetGreatDim}
%<*co:NoLiSetGreatDim>
No linearly independent set can have a size greater than the dimension of the
enclosing space.
%</co:NoLiSetGreatDim>
\end{corollary}

% B40.dimension-bound proof; src/vs/vs3.tex; source lines 1825-1831
\begin{proof}
%<*pf:NoLiSetGreatDim>
The proof of \nearbytheorem{th:AllBasesSameSize} 
never uses that \( D \) spans the space,
only that it is linearly independent.
%</pf:NoLiSetGreatDim>
\end{proof}

% B40.basis-extension statement; src/vs/vs3.tex; source lines 1857-1861
\begin{corollary} \label{cor:LIExpBas}
%<*co:LIExpBas>
Any linearly independent set can be expanded to make a basis.
%</co:LIExpBas>
\end{corollary}

% B40.basis-extension proof; src/vs/vs3.tex; source lines 1863-1874
\begin{proof}
%<*pf:LIExpBas>
If a linearly independent set 
is not already a basis then it must not span the space.
Adding to the set a vector that is not in the span 
will preserve linear independence by 
Lemma II.\ref{lm:AddVecLiSetIsLiIffVecNotInSpan}.
Keep adding until the resulting set does span the space, 
which the prior corollary
shows will happen after only a finite number of steps.
%</pf:LIExpBas>
\end{proof}
